\chapter{A - The work flow}
In this chapter I will recount the work I have done and the steps taken to achieve the final goal.
\section{The travelling waves}
Before one can start solving travelling KdV waves in time one needs to build travelling KdV waves. 
For me this started with looking at the Airy equation and implementing the $sech^2$ wave, both as a univariale, and a bivariate function. 
From this I worked on methods of solving them forward in time, and comparing them to their analytic counterparts, however that is a topic for the next chapter.
The form of the Airy functions intreged me and I tried to no avail to create my own, the method was unstable numerically and I consiged and used the prebuilt functions from scipy.
From this point on, I started looking at newton solvers for the travelling waves. 

    \subsection{Naive Newton}
    Wether the name has any theoretical annotations I do not know, the method I created however was quite naive in its implemntation compared to later once. 
    To start of with I i worked on the physical side, and implemented a newton solver - without jacobian - and found I could create waves of small amplitude, and periods close to $2 \pi$. 
    These waves were to the naked eye simmilar to the sine waves but solved the KdV equation and had strict numerical restrictions due to the third derivative that needed to be approximated. 
    These restrictions made me cast a wider net, and I found that the Newton solver also had a Fourier counterpart

    \subsection{Spectral Newton}
    In the spectral space the equations became sparse, the travelling waves do not contain a wide spectrum, and is real and symmetric.
    The equations themselves are not that different and teh change was implemented quite quickly. 
    As it would turn out however that the main issue of beeing limited in both amplitude and period persited, and another avenue had to be explored. 
    This came by ways of continuation, by incrementing the amplitude up to the desired size older stable solutions 
    were built oppon and made larger amplitudes possible. That is one sovles the problem for a small amplitude
    and use that solution as the initial data for a larger amplitude, until the given amplitude is reached,
    given it converges. 




This concludes the introduction to the travelling waves and the newton solvers. 


\section{Time-dependent KdV}
Returning to the prospect of solving PDE's I started by looking into linear problems, and ommiting the nonlinearities. 
Reading [Trefethen et al.] and some of the contempory methods for solving the full equation, 
in addition to my own experience with the heat equation and the heat kernel, I found the integrating factor mehod appeling.
By solving the equation in this way I "dicovered" that all linear PDE's of this type could be solved by convolution with a kernel. 
The form of the kernel was given by the linear relation in question. Testing it on the transport equation and heat equation gave the expected results,
and I was convinced: The integrating factor method was hot shit. 

My dissilutions were eventually cleared up, but not before I created a full solver using both:
integrating factor and split-step to work on the linear and nonlinear parts separatly. 
If I recall correctly the split step failed when I cleared up in a dealiasing issue,
now that I have removed the dealiasing I could warrent a revisit. (adding it to the todo)
To start of with I, again, looked at the $sech^2$ waves and compared my solution to the analytical one.
From here I worked on dealiasing of the nonlinearity - at the time I thought this would be a problem. And I compared both padding and masking, but found no 
difference in the solver, as masking was the simplest one to implement I continued using this. 

New to dealiasing and this \textit{complicated} sovlers, I turned to AI (Claude) for help in how to apply it. 
The answer given seemed a bit off at the time, but I did'nt know better and added the masking both inside the nonlinearity and outside.
Solving it this way my split-step and integrating-factor both worked and were of order 2. 

After reading more theory and asking more pointed questions to both my supervisors and AI I changed the way my masking was applied,
doing this did not change the results of the split-step method but my IF one failed. Given that creating solvers for the 
PDE was only a small step in the thesis, I moved on with the SS method. 

\section{Recovery on KdV data}
With both a travelling wave generator and a forward solver, I could now start on the recovery process. 
To get a better feel for the processes I made a solver that only tried to recover the linear symbol. 
Doing it this way ment I could make one wave and then use LSM on only the unknown part of it.
In my first attempt I used the travelling wave assumption and made it work on a static wave profile.
The recovered symbol was accurate out to $k = 11$, where - using the assumptions from Theory - the strength of the signal approched the tolerance of the "wave generator".
From this platform adding the nonlinearity was a simple procedure, as all I had to do was rework the matrix multiplications. 


\section{Full dispersion}
After making speciffic solvers for the KdV equation, as described above. I set out to generalize, trying to capture the respective relations for 
both Whithams m and the fractional KdV (fKdV). This led me down the path of objective programming. I stared working on a general class that could be 
initiated as any one of the model problems I wanted to look at, that is I created an abstract sovler class and made a child for each of the equations. 
Doing it this way I could easily set up test cases and plot the important graphs without having to rewrite or copy code. 
But the set up phase took some time. 

In addition to looking into different equations, I also looked into how to represent the equation in my LSM solver.
The main point of contension here was the time derivative. In the article by my supervisors, they estimated it using the Fourier transform in time 
and only used a few (3 iirc) datapoints in time, though this was also limited by the position of the soliton in the frame. 
In my own experiments having to refactor the time data so that it is periodic was an issue I didnt find a good solution to. 
I made individual times for each soliton such that they would move exactly one period in time. This solution however only works 
on travelling waves where the motion can be periodic in time. Looking at eg. a Gaussian blob this method can not be used. 